% ============================================================
%  Physics-Based PV Generation Prediction for IoT SCADA
%  A Business R&D Article
%  Compile with: pdflatex + bibtex (or lualatex / xelatex)
% ============================================================

\documentclass[11pt, a4paper]{article}

% ---- Encoding & font ----------------------------------------
\usepackage[T1]{fontenc}
\usepackage[utf8]{inputenc}
\usepackage{palatino}           % Palatino body font (LinkedIn-ish warmth)
\usepackage{microtype}          % Improved kerning & protrusion

% ---- Math & symbols ----------------------------------------
\usepackage{amsmath}            % \text{} inside math mode
\usepackage{amssymb}            % \checkmark, standard math symbols
\usepackage{eurosym}            % \EUR{} euro sign

% ---- Page geometry ------------------------------------------
\usepackage[
  top=2.5cm, bottom=2.5cm,
  left=2.8cm, right=2.8cm,
  headheight=14pt              % silence fancyhdr \headheight warning
]{geometry}

% ---- Colours ------------------------------------------------
\usepackage[dvipsnames, table]{xcolor}
\definecolor{brandblue}{HTML}{0A66C2}   % LinkedIn brand blue
\definecolor{brandgrey}{HTML}{F3F2EF}   % LinkedIn light grey background
\definecolor{accentorange}{HTML}{E07B39}
\definecolor{accentgreen}{HTML}{3B7D3A}
\definecolor{lightgrey}{HTML}{6B7280}
\definecolor{darkgrey}{HTML}{1F2937}

% ---- Section styling ----------------------------------------
\usepackage{titlesec}
\titleformat{\section}
  {\large\bfseries\color{brandblue}}
  {}
  {0em}
  {}
  [\vspace{-0.4em}\rule{\linewidth}{1.5pt}\vspace{0.2em}]

\titleformat{\subsection}
  {\normalsize\bfseries\color{darkgrey}}
  {}
  {0em}
  {}

\titleformat{\subsubsection}
  {\small\bfseries\color{lightgrey}\itshape}
  {}
  {0em}
  {}

% ---- Paragraph spacing --------------------------------------
\usepackage{parskip}
\setlength{\parskip}{0.6em}
\setlength{\parindent}{0pt}

% ---- Lists --------------------------------------------------
\usepackage{enumitem}
\setlist[itemize]{
  leftmargin=1.4em,
  topsep=0.2em,
  itemsep=0.2em,
  label={\color{brandblue}\textbullet}
}
\setlist[enumerate]{
  leftmargin=1.6em,
  topsep=0.2em,
  itemsep=0.2em
}

% ---- Tables ------------------------------------------------
\usepackage{booktabs}
\usepackage{tabularx}
\usepackage{array}
\usepackage{multirow}
\renewcommand{\arraystretch}{1.3}

% ---- Callout / highlight boxes ----------------------------
\usepackage{tcolorbox}
\tcbuselibrary{skins, breakable}

% Insight box (blue tint)
\newtcolorbox{insightbox}[1][]{
  enhanced,
  breakable,
  colback=brandblue!6,
  colframe=brandblue,
  fonttitle=\bfseries\small,
  title={#1},
  left=6pt, right=6pt, top=4pt, bottom=4pt,
  arc=3pt
}

% Warning / gap box (orange tint)
\newtcolorbox{gapbox}[1][]{
  enhanced,
  breakable,
  colback=accentorange!8,
  colframe=accentorange,
  fonttitle=\bfseries\small,
  title={#1},
  left=6pt, right=6pt, top=4pt, bottom=4pt,
  arc=3pt
}

% Stat highlight (green)
\newtcolorbox{statbox}{
  enhanced,
  colback=accentgreen!7,
  colframe=accentgreen,
  left=6pt, right=6pt, top=3pt, bottom=3pt,
  arc=3pt
}

% ---- Hyperlinks --------------------------------------------
\usepackage[
  colorlinks=true,
  linkcolor=brandblue,
  citecolor=brandblue,
  urlcolor=brandblue
]{hyperref}

% ---- Graphics & floats ------------------------------------
\usepackage{graphicx}
\usepackage{float}

% ---- Header / footer --------------------------------------
\usepackage{fancyhdr}
\pagestyle{fancy}
\fancyhf{}
\fancyhead[L]{\small\color{lightgrey}\textit{ThingsNode · S.\,M.\,Pasindu Upeksha Senevirathne}}
\fancyhead[R]{\small\color{lightgrey}March 2026}
\fancyfoot[C]{\small\color{lightgrey}\thepage}
\renewcommand{\headrulewidth}{0.3pt}
\renewcommand{\headrule}{\color{brandblue}\hrule width\headwidth height\headrulewidth}

% ---- Bibliography ------------------------------------------
\usepackage[
  backend=bibtex,
  style=numeric,
  sorting=none,
  maxbibnames=3
]{biblatex}
\addbibresource{article.bib}

% ---- Utility macros -----------------------------------------
\newcommand{\highlight}[1]{\textbf{\color{brandblue}#1}}
\newcommand{\kw}[1]{\textit{#1}}
\newcommand{\metric}[2]{\textbf{#1}\,\textcolor{lightgrey}{\small(#2)}}

% ============================================================
\begin{document}

% ---- Title block -------------------------------------------
\begin{center}
  {\fontsize{22}{28}\selectfont\bfseries\color{darkgrey}
    Physics-Based Solar PV Generation Prediction\\[4pt]
    \color{brandblue}for IoT SCADA Systems}\\[10pt]
  {\large\color{lightgrey}
    A Practical R\&D Journey: Physics-Based PV Modelling on AWS Lambda,\\
    from PVsyst Parameters to a Production IoT Pipeline}\\[14pt]
  \rule{0.6\linewidth}{0.6pt}\\[8pt]
  {\small\color{lightgrey}
    S.\,M.\,Pasindu Upeksha Senevirathne \quad|\quad ThingsNode \quad|\quad March 2026 \quad|\quad LinkedIn}
\end{center}

\vspace{0.6em}

% ---- Author note -------------------------------------------
\begin{insightbox}[About This Article]
This article documents an engineering project I built during my internship
at ThingsNode: a physics-based PV generation prediction system for a 55\,kW
multi-orientation rooftop solar installation in Sri Lanka, deployed on an
AWS IoT pipeline. It covers the scientific methods chosen, the approaches
explored and discarded, and the accuracy achieved in real operational conditions.
\end{insightbox}

\vspace{0.5em}

% ============================================================
\section{The Problem Nobody Has Fully Solved}
% ============================================================

Monitoring a solar PV plant in real time means answering one fundamental
question at every telemetry tick: \highlight{is this plant performing as expected?}

To answer that, you need two numbers: \emph{actual generation} (from the meter)
and \emph{expected generation} (from a model). The first is easy. The second is
surprisingly hard to do well, cheaply, and fast enough for real-time SCADA.

The core tension is this:

\begin{itemize}
  \item \textbf{Physics models} (pvlib, PVsyst) are accurate, but running
    them inside an IoT pipeline --- triggered at every telemetry message, at
    fleet scale --- requires an architecture that most teams don't have.
  \item \textbf{Commercial cloud APIs} (SolarEdge, AlsoEnergy PowerTrack,
    Sungrow iSolarCloud) are tightly coupled to their own hardware and
    subscription plans, and do not work for third-party inverters or
    multi-vendor fleets.
  \item \textbf{Simple regression models} trained on historical data require
    ground-truth generation labels, take weeks to calibrate per plant, and
    break the moment a module is added or shading changes.
  \item \textbf{Nameplate-based estimators} (``GHI × rated efficiency × area'')
    ignore orientation, temperature, clipping, and IAM --- error margins of
    20--30\% are common.
\end{itemize}

We set out to close this gap for a real operational deployment:
a \textbf{55\,kW multi-orientation rooftop PV system} on a school building
in North-Central Sri Lanka, connected to an AWS IoT pipeline and a Meteocontrol
data logger, with 8 differently-oriented roof planes spanning all compass
quadrants.

The stakes extend beyond individual asset economics. In early 2025, Sri Lanka
experienced grid instability linked to the rapid, unpredicted surge in rooftop
solar generation on a low-demand day --- an event that highlighted the systemic
risk of not knowing, in real time, how much generation a distributed PV fleet
is actually producing vs what was expected. Accurate, per-site expected
generation is not only an O\&M tool; it is a grid-stability input.

At plant level, the impact is equally concrete. A persistent 10\% soiling
loss on a 55\,kW plant means roughly 30--36\,kWh of missed generation every
clear day --- \$1{,}500--2{,}500 per year in avoided-cost or feed-in revenue,
continuously invisible without a physics baseline to compare against.

% ============================================================
\section{The Market Landscape: What Others Do (and Don't)}
% ============================================================

\subsection{Commercial Monitoring Platforms}

Most solar monitoring vendors are vertically integrated: they sell the inverter,
the logger, and the dashboard as a bundle. Expected generation is typically
backed by a simple clearness-index model or a manufacturer-supplied yield
curve --- neither is PVsyst-consistent or orientation-aware.

\begin{table}[H]
\centering
\caption{Market solutions vs our approach}\label{tab:market}
\small
\begin{tabularx}{\linewidth}{>{\bfseries}p{2.6cm} X X X}
\toprule
Feature &
\color{brandblue}ThingsNode Engine &
Commercial Platforms &
Open-Source (pvlib) \\
\midrule
Vendor hardware lock       & None         & Often required     & None \\
Serverless cloud (AWS Lambda)& \color{accentgreen}\checkmark & Proprietary cloud  & \color{accentorange}$\times$ \\
PVsyst-consistent physics  & \color{accentgreen}\checkmark & Partial            & \color{accentgreen}\checkmark \\
Multi-orientation support  & \color{accentgreen}\checkmark & Per-inverter only  & \color{accentgreen}\checkmark \\
Free irradiance sources    & NASA/Open-Meteo & Proprietary       & Manual \\
No vendor hardware lock    & \color{accentgreen}\checkmark & Often required     & \color{accentgreen}\checkmark \\
Multi-source data fallback & \color{accentgreen}\checkmark & \color{accentorange}$\times$ & Manual \\
Per-plant calibration      & \color{accentgreen}\checkmark & Black-box          & Manual \\
Open source                & Planned       & \color{accentorange}$\times$ & \color{accentgreen}\checkmark \\
\bottomrule
\end{tabularx}
\end{table}

\subsection{What the Market Does Well (and We Should Acknowledge)}

\begin{itemize}
  \item \textbf{SolarEdge, SMA, Huawei}: Deep hardware integration gives them
    precise DC string-level diagnostics we cannot match from irradiance alone.
  \item \textbf{Solcast}: Class-leading satellite irradiance with 15-minute
    resolution and uncertainty bands --- we use Solcast as a data source, not a
    competitor.
  \item \textbf{Aurora Solar, PVCase}: Superior shading engine using 3-D models
    and LiDAR --- far beyond our scope.
  \item \textbf{PVsyst}: Industry gold standard for bankable energy yield
    assessments. Our model is designed to be \emph{consistent} with PVsyst,
    not to replace it.
\end{itemize}

\subsection{What the Market Does Not Do Well}

\begin{itemize}
  \item \textbf{Multi-vendor SCADA at accessible cost}: Enterprise APM platforms
    (AlsoEnergy, Power Factors) \emph{can} handle heterogeneous fleets, but
    they are priced for utility-scale operators. The barrier for independent
    C\&I operators is not technical --- it is cost and proprietary integration
    complexity. A self-hosted physics engine removes both.
  \item \textbf{Open API access}: Major OEMs (e.g.\ SolarEdge) have
    progressively restricted third-party API access, forcing independent
    analytics platforms onto rate-limited ``associated account'' channels.
    Relying on standard data logger telemetry (Meteocontrol, SMA Data Manager)
    rather than cloud-to-cloud OEM APIs sidesteps this dependency entirely.
  \item \textbf{Edge / offline expected generation}: A rule chain firing
    every 15 minutes across thousands of devices cannot make an external
    API call per tick without latency and cost implications.
  \item \textbf{Transparent, auditable physics}: Most commercial expected
    generation values are black-box numbers. Our model is fully inspectable,
    config-driven, and PVsyst-traceable.
  \item \textbf{Cost inefficiency for independent operators}: Physics-grade
    PV modelling typically requires a PVsyst license
    (\textasciitilde\EUR{}2{,}000/year) or a commercial monitoring
    subscription at \$10--\$50 per kW per year. For most operators, this
    cost buys a black-box number with no way to audit, debate, or correct it.
\end{itemize}

% ============================================================
\section{Our System --- What We Have (and Why)}
% ============================================================

\subsection{The Physics Engine}

The core of the system is a full, PVsyst-consistent physics chain implemented
in Python using pvlib \cite{holmgren2018pvlib}. It is the engineering
source of truth for the entire pipeline.

\textbf{The computation chain, in order:}

\begin{enumerate}
  \item \textbf{Irradiance ingestion} --- from multiple satellite and reanalysis
    sources: Solcast, NASA POWER, Open-Meteo, ERA5, or local measurement files
  \item \textbf{Solar position} --- NREL Solar Position Algorithm (sub-0.01°
    accuracy \cite{reda2004solar})
  \item \textbf{POA irradiance transposition} --- Perez diffuse model
    \cite{perez1990set} applied independently to each of the 8 roof orientations
  \item \textbf{Far shading} --- a multiplicative scalar factor from the plant's
    PVsyst site report
  \item \textbf{AOI \& IAM correction} --- angle-of-incidence and incidence
    angle modifier table, extracted from PVsyst and interpolated per orientation
    \cite{deline2011broken}
  \item \textbf{Cell temperature} --- Sandia Array Performance Model (SAPM)
    thermal model \cite{king2004photovoltaic}, driven by POA irradiance and
    wind speed
  \item \textbf{DC power} --- temperature-corrected PVWatts-style formulation:
    $P_{DC} = \text{POA}_{\text{opt}} \times \eta_{STC}
    \times (1 + \gamma_P (T_{cell} - 25)) \times k_{DC}$
  \item \textbf{DC loss chain} --- 5 sequential multipliers from PVsyst loss
    diagram:
    soiling (3.0\%) $\to$ LID (1.4\%) $\to$ module quality ($-$0.8\% gain)
    $\to$ mismatch (1.7\%) $\to$ DC wiring (0.9\%)
  \item \textbf{Load-dependent inverter efficiency} --- 21-point efficiency
    curve from PVsyst, linearly interpolated at runtime
  \item \textbf{Plant-level clipping} --- applied once at 55\,kW AC after
    summing all 8 orientations
  \item \textbf{AC wiring loss} --- 0.3\% applied to the clipped plant output
\end{enumerate}

\begin{insightbox}[Why we chose these methods]
Each step was specifically chosen for PVsyst consistency. The Perez model
outperforms simpler isotropic diffuse models by 5--10\% RMSE on tilted
multi-orientation arrays \cite{gueymard2009direct}. SAPM's inclusion of
wind speed is critical in tropical coastal climates: research on tropical
PV sites shows SAPM cell temperature errors of 2--5\,K, compared to up
to 14\,K for static NOCT models \cite{king2004photovoltaic}. At
$\gamma_P = -0.34\%/^\circ$C, each 3\,K error propagates to a 1\% DC
power bias. The load-dependent inverter curve eliminates a systematic 1--3\%
over-prediction at partial load that a constant-efficiency assumption introduces.
\end{insightbox}

\subsection{The Irradiance Fallback Architecture}

No single free irradiance source is reliable in all weather conditions.
The prediction system draws simultaneously from four independent sources
--- Solcast, NASA POWER, Open-Meteo, and ERA5 --- and assembles each
weather field (GHI, DNI, DHI, air temperature, wind speed) from the
best available source per timestamp:

\[
\text{Solcast} \to \text{NASA POWER} \to \text{Open-Meteo} \to \text{ERA5}
\]

The cascade priority reflects empirical accuracy: 2025 global benchmarking
shows Solcast achieving a GHI Mean Absolute Percentage Error (MAPE) of
\textasciitilde{}10.7\% vs ERA5's \textasciitilde{}18.4\% across diverse
climates \cite{ineichen2011global}. Both converge during deep convective
monsoon events --- the fundamental limit of satellite and reanalysis products.
ERA5 provides temperature and wind but not irradiance, so it serves as a
weather-only fallback.

\subsection{The Explored Path: A Calibrated Approximation for Edge Deployment}

Before committing to a cloud architecture, we investigated whether a physics
approximation could run directly inside an IoT edge runtime --- specifically
a ThingsBoard rule chain --- to eliminate any cloud round-trip entirely.
This required translating the full physics model into a self-contained
JavaScript approximation, then calibrating it against the reference model.

The calibration strategy treated the Python physics engine as a
\kw{physics oracle}: 2.5~million synthetic weather--power samples spanning
GHI 0--1{,}200\,W/m$^2$, temperature 15--45\,$^\circ$C, and wind 0.1--15\,m/s
were generated and used to minimise the approximation error via global
differential evolution optimisation \cite{storn1997differential}.

\begin{statbox}
\centering
\textbf{Approximation accuracy:} RMSE = 1.43\,kW \quad MAE = 0.67\,kW \quad
R$^2$ = 0.9923 \quad Energy error = 1.49\% \quad (validated on 2.5M samples)
\end{statbox}

\textbf{Why this path was rejected for production:} The approximation achieves
high correlation with the physics reference internally, but introduces systematic
bias from a simplified diffuse irradiance model and flat inverter efficiency.
These biases, small per hour, compound into meaningful errors over monthly
reporting cycles. A 200--500\,ms serverless invocation is imperceptible on a
15-minute telemetry cadence, so the latency advantage offered no practical
benefit. Accuracy was chosen over speed.

\subsection{The Production Architecture: Physics on AWS Lambda}\label{sec:aws}

The production system runs the full Python physics engine as a serverless
AWS Lambda function, triggered by AWS IoT Core rule actions on each
incoming telemetry event.

\textbf{Data flow:}
\begin{enumerate}
  \item AWS IoT Core receives 15-minute telemetry from the plant's Meteocontrol
    data logger
  \item An IoT rule routes each message to a Lambda function
  \item Lambda fetches latest irradiance from the best available source,
    runs the full physics engine, and computes expected AC power
  \item The expected generation value is published back to the IoT topic
    and surfaced on the monitoring dashboard
\end{enumerate}

This preserves the complete Perez POA transposition, SAPM thermal model,
load-dependent inverter efficiency curve, and 5-factor DC loss chain ---
with no accuracy compromise from approximation.

\subsection{The Linear Regression Surrogate (Archived)}

A much simpler alternative was also explored: a 5-feature polynomial
regression of the form
$P_{AC} = a_0 + a_1 GHI + a_2 GHI^2 + a_3 DNI
+ a_4 (GHI \cdot \Delta T) + a_5 (GHI \cdot v_{wind})$,
fit to one year of physics-model outputs. The model runs nearly instantaneously
but achieves R$^2$ = 0.42 and MAPE = 130\% on held-out hourly data. Monthly
energy error is only 1.5\% --- but this is a statistical artefact: the
model over-predicts in the morning and under-predicts in the afternoon in
ways that cancel over 24 hours, not because the physics is correct.
It is archived for contexts where only daily energy totals, not hourly
profiles, are needed.

% ============================================================
\section{Our System --- What We Don't Have (and Why)}
% ============================================================

\begin{gapbox}[Intentional scope limits]
The gaps below are not oversights. Each was evaluated and either deferred for
good engineering reasons, or is simply out of project scope by design.
\end{gapbox}

\subsection{No Near-Shading Module}
We model \kw{far shading} (a scalar horizon factor from PVsyst) but not
\kw{near shading} --- row-to-row, parapet, or chimney shading that varies hour
by hour. This requires a 3-D geometric engine (as in PVsyst or Aurora Solar)
and per-plant 3-D measurement. Our target deployment is a single-family or
commercial rooftop where PVsyst has already been used at commissioning; we
inherit PVsyst's shading output rather than recompute it.

\subsection{No Spectral Correction}
PVsyst applies a module-specific spectral response correction that adjusts
yield by up to 1--2\% depending on air mass and aerosol optical depth
\cite{wright2020spectral}. We do not implement this. At low altitude (88\,m)
and tropical latitudes, the effect is small and within our target accuracy band.

\subsection{No Annual Degradation}
Module output degrades at roughly 0.5--0.8\% per year \cite{jordan2016robust}.
We apply LID (first-year degradation) but not an ongoing annual degradation rate
($R_d$). For near-real-time anomaly detection on a recently commissioned plant,
seasonal irradiance variation dominates this signal. However, there is an
important long-term caveat: without progressively adjusting the expected
generation baseline to account for cumulative degradation, a multi-year PR trend
will artificially register as worsening equipment performance even when the plant
is behaving exactly as expected. Degradation tracking is planned as a future
module and is necessary to prevent false-positive fault alerts after year 3--5.

\subsection{No Battery / Self-Consumption Offset}
The model predicts gross AC generation. If the plant has a battery storage
system or significant self-consumption, the site meter reading will differ from
generation. As co-located Battery Energy Storage Systems (BESS) become
standard in new C\&I solar installations globally, this is a growing scope
boundary: reconciling predicted gross PV generation against a net-metered
site export requires modelling battery dispatch, state-of-charge, and
peak-shaving algorithms. This is architecturally out of scope for the
current generation prediction engine but is an identified future extension.

\subsection{No Load-Dependent Inverter Curve in the JS Engine}
The JavaScript physics engine uses a flat 98\% inverter efficiency. Implementing
a full 21-point interpolation in Nashorn JavaScript is possible but
significantly increases the rule chain computation budget. At $<$3\% energy
error, the flat efficiency is acceptable for real-time SCADA use. The Python
reference model does use the full curve.

\subsection{No Self-Calibrating Loss Factors}
Our loss factors (soiling, LID, mismatch) are extracted from PVsyst at
commissioning and remain static. A more advanced system would recalibrate
soiling losses automatically as actual vs expected divergence grows. This
requires several months of high-quality actuals and a calibration algorithm ---
planned but not implemented.

\subsection{No PVsyst Annual Validation}
We have not yet run a full PVsyst simulation for this plant and compared
annual energy totals. This is the standard bankable-yield validation step, and
it remains pending. Our model is designed to be PVsyst-consistent, but
consistency is assumption-based until formally validated.

% ============================================================
\section{Technical Deep Dive: Key Engineering Decisions}
% ============================================================

\subsection{Why Multi-Orientation Kills Simple Regression}

With a single south-facing array, GHI maps reasonably to POA ($\text{POA}
\approx \text{GHI} \times f(\text{tilt})$). With 8 orientations spanning SE,
NNW, ENE, and ESE simultaneously, the same GHI value at 08:00 and at 14:00
produces completely different total POA because east-facing arrays peak in
the morning and west-facing arrays peak in the afternoon. A regression on GHI
cannot distinguish these --- our linear surrogate's R$^2$ = 0.42 proves this
empirically.

\begin{insightbox}[The multi-orientation insight]
The 8-orientation plant spreads energy production over a longer daily window,
reducing clipping and smoothing output. On this plant, east-facing arrays
contribute approximately 35\% of daily energy before noon; west-facing arrays
contribute approximately 30\% after noon --- with the south-facing sections
bridging the midday peak. Any prediction model that doesn't resolve each
orientation's shifting POA independently will fail on mixed-azimuth arrays.
\end{insightbox}

\subsection{Plant-Level vs Per-Orientation Clipping}

Early code clipped each orientation's AC power independently before aggregation.
This is mathematically incorrect. A multi-MPPT string inverter clips the
\emph{sum} of all DC inputs at the plant AC rating, not each sub-array
independently. Per-orientation clipping predicted approximately 3\% more
clipping loss than reality --- it penalised sub-arrays whose individual DC
peak is below the threshold, even when their combined DC total exceeds it.
Summing all orientation DC outputs first, then applying the single inverter
clipping threshold, accurately mirrors the electrical behaviour of the physical
hardware. We corrected this before any benchmarks were run.

\subsection{Why Differential Evolution, Not Gradient Descent}

During the JS edge exploration phase, we needed to tune 8 non-linear model
parameters (SAPM thermal constants and Perez diffuse weights) simultaneously.
SAPM and Perez terms interact --- changing the thermal model shifts cell
temperature, which shifts DC power, which changes the load point on the inverter
curve. This creates a highly non-convex loss surface. Gradient-based methods
(Newton--Raphson, L-BFGS) require derivative information and are acutely
sensitive to starting conditions on such surfaces --- in early tests they
converged reliably only when manually initialised near the correct region, which
is impractical for a new plant. Differential evolution
\cite{storn1997differential}, a population-based global optimiser, explores the
full parameter space simultaneously and converged to the solution in 112
iterations and 13{,}677 evaluations with no starting-point tuning.

\subsection{Synthetic Training Data Generation}

Rather than using real historical weather data (which is limited to the period
available from NASA POWER), we generated 2.5 million random, physically-plausible
weather samples covering GHI 0--1200\,W/m$^2$, temperature 15--45$^\circ$C,
and wind 0.1--15\,m/s across all daylight hours of a full tropical year. This
``physics oracle'' approach is similar to sim-to-real transfer in robotics
\cite{tobin2017domain}: use a high-fidelity simulator to generate copious
training data that empirical observation cannot.

% ============================================================
\section{Performance \& Accuracy Results}
% ============================================================

\subsection{JS Physics vs Python Reference (Internal Accuracy)}

\begin{table}[H]
\centering
\caption{JS Physics engine vs Python pvlib reference}
\label{tab:jsvsphysics}
\small
\begin{tabularx}{\linewidth}{l >{\centering\arraybackslash}X >{\centering\arraybackslash}X}
\toprule
Metric & Value & Interpretation \\
\midrule
RMSE       & 1.46\,kW  & 2.65\% of 55\,kW rating \\
MAE        & 0.68\,kW  & 1.2\% of rating \\
R$^2$      & 0.9921    & Near-perfect correlation \\
Energy error & 1.46\%  & Monthly totals match well \\
MAPE       & 7.76\%    & Elevated at dawn/dusk (low power) \\
Speed      & 0.026\,ms & 28$\times$ faster than Python \\
Validated on & 2.5M samples & Synthetic, full input-space coverage \\
\bottomrule
\end{tabularx}
\end{table}

\subsection{Model Comparison on Real 2024 Test Data}

Using 1{,}357 held-out hourly samples from 2024 NASA POWER data (73 calendar days,
never seen during calibration):

\begin{table}[H]
\centering
\caption{Model comparison on real held-out test set (JS vs Python model output)}
\label{tab:testset}
\small
\begin{tabularx}{\linewidth}{l *{5}{>{\centering\arraybackslash}X}}
\toprule
Model & R$^2$ & MAE (kW) & MAPE & Energy error & Speed \\
\midrule
JS Physics  & 0.969 & 1.34 & 20.1\% & +5.47\% & 0.026\,ms \\
JS Surrogate & 0.420 & 9.44 & 129.7\% & +1.50\% & 0.011\,ms \\
\bottomrule
\end{tabularx}
\end{table}

\begin{gapbox}[Why the surrogate's energy accuracy is misleading]
The surrogate achieves only 1.5\% energy error but R$^2$ = 0.42. This is
because the model systematically over-predicts at low irradiance and
under-predicts at high irradiance in ways that \emph{happen to cancel} over a
full day. For hourly SCADA performance monitoring, it is useless. For monthly
billing reconciliation, it might be acceptable --- but only if you accept that
you cannot detect event-level anomalies.
\end{gapbox}

\subsection{Against Real Meteocontrol Actuals (December 2025)}

In the December 2025 benchmark window, available Meteocontrol actuals in the
repository overlap on only two days (11 and 12 December 2025). On that overlap,
all three models over-predicted generation. For example: on 12 December 2025,
actual generation was \textbf{76\,kWh}; the Python model prediction was
\textbf{263\,kWh} --- a 246\% over-prediction.

Python +15.8\%, JS Physics +27.0\%, JS Surrogate +18.3\%.
These numbers reflect the irradiance error, not the physics error.

\begin{insightbox}[Why this is a data problem, not a physics problem]
December in Sri Lanka is the North-East monsoon season. In December 2025,
Cyclone Ditwah developed in the Bay of Bengal and subjected Sri Lanka and
Southern India to intense cloud cover and rainfall, slashing irradiance
significantly below seasonal norms. Solcast's own published irradiance
analysis for this period confirms the anomaly \cite{solcast2025cyclone}.
All reanalysis (NASA POWER, ERA5) and satellite products (Solcast) universally
struggle to resolve micro-scale cloud optical depth during cyclonic events --
Solcast reports GHI hourly MAPE exceeding 30\% in tropical convective
regimes during monsoon.
This is not a model failure --- it is the fundamental limitation of irradiance
determination under deep convective cloud cover. No physics model can produce
the right answer when given wrong irradiance.
\end{insightbox}

\subsection{Irradiance Source Comparison (10-day window)}

\begin{table}[H]
\centering
\caption{Daily prediction accuracy by irradiance source (December 2025, 10 days vs Meteocontrol)}
\label{tab:sources}
\small
\begin{tabular}{l r r r r}
\toprule
Source & Daily MAE (kWh) & Bias (kWh) & R$^2$ & Availability \\
\midrule
NASA POWER  & 96.2 & +32.7 & $-$0.68 & Free, global \\
Solcast     & 97.4 & +33.9 & $-$0.71 & Paid / free tier \\
Open-Meteo  & 111.4 & +47.9 & $-$1.55 & Free, no key \\
ERA5        & 96.8 & +33.3 & $-$0.70 & Free, delayed \\
\bottomrule
\end{tabular}
\end{table}

Negative R$^2$ across all sources during this window reflects the monsoon
outlier problem. NASA POWER and ERA5 perform nearly identically (both are
GEOS-derived reanalysis products). Open-Meteo shows higher bias during cloudy
periods.

% ============================================================
\section{What the Research Says}
% ============================================================

\subsection{Physics vs Machine Learning for PV Prediction}

The literature is evolving. Diagne \textit{et al.} \cite{diagne2013review}
showed that physical models consistently outperform statistical models beyond
the 6-hour horizon because they do not require the future to resemble the
training period. Mellit \textit{et al.} \cite{mellit2010artificial} demonstrated
that ANN-based models achieve competitive short-term (1-hour) accuracy
\emph{but only} when trained on several years of site-specific data --- a
condition rarely met for newly commissioned plants.

By 2025--2026, the industry has largely moved toward Physics-Informed Machine
Learning (PIML) hybrids: neural networks whose loss functions encode physical
constraints (solar geometry, thermal exchange) to prevent outputs that violate
thermodynamic laws. These hybrids reduce MAE by 15--30\% compared to standalone
models on short-term tasks. However, for \emph{auditable operational baselines}
--- where an O\&M engineer must trace every prediction to a physical assumption
and debate it --- pure physics chains retain a decisive advantage: they require
no labelled historical data, generalise immediately to new plants from first
principles, and every intermediate output is interpretable. Our empirical result
(surrogate R$^2$ = 0.42 vs JS Physics R$^2$ = 0.97) confirms this in a
production setting with only one year of training data.

\subsection{Irradiance Transposition}

The Perez all-weather model \cite{perez1990set} remains the benchmark for
tilted-surface irradiance decomposition. Gueymard \cite{gueymard2009direct}
evaluated 54 transposition models and found Perez outperformed isotropic, Hay,
and Reindl models by 5--10\% RMSE on non-south-facing arrays. This directly
motivated our choice over the simpler isotropic model in the surrogate.

\subsection{Cell Temperature Modelling}

King \textit{et al.} \cite{king2004photovoltaic} showed that the SAPM thermal
model (including wind speed) explains $>$95\% of cell temperature variance for
rack-mounted modules. The NOCT-based model, still widely used in commercial
platforms, ignores wind and introduces systematic 3--5\,K bias in coastal sites.
At $\gamma_P = -0.34\%/^\circ$C, a 4\,K error in cell temperature translates
to approximately 1.4\% DC power error.

\subsection{Multi-Orientation Complex Rooftops}

Haghdadi \textit{et al.} \cite{haghdadi2017method} demonstrated that
orientation-disaggregated modelling is necessary for complex rooftops:
aggregate models (single representative tilt/azimuth) under-predict energy
by 5--15\% on mixed-orientation arrays. Our per-orientation loop with
area-weighted aggregation addresses this finding directly.

\subsection{IoT Platform Constraints}

Performance benchmarks for ThingsBoard rule chains \cite{thingsboard2024perf}
show that a single script node processing 10{,}000 messages per second requires
sub-0.1\,ms execution. Our JS engine at 0.026\,ms comfortably meets this
requirement. The Python pvlib stack at 0.74\,ms per call would fail at fleet
scale ($>$1{,}000 sites at 15-minute intervals).

\subsection{Sim-to-Real Transfer for Hybrid Models}

Our calibration approach --- generating synthetic data from a physics model to
train a faster approximation --- has precedent in robotics (domain randomisation
\cite{tobin2017domain}) and PV modelling. Inman \textit{et al.}
\cite{inman2013solar} used physics-model outputs as labels to train
computationally efficient surrogates, reporting that physics-informed training
sets produce surrogates that generalise better than purely data-driven ones.
The simple polynomial regression we built demonstrates the floor of this
approach. Future surrogates would benefit from gradient-boosted ensemble
methods (e.g.\ XGBoost or LightGBM), which handle the non-linear, interaction-heavy
relationship between irradiance, temperature, wind, and multi-orientation POA
far more effectively than a polynomial basis, without sacrificing
the physics-oracle training strategy.

% ============================================================
\section{Competitive Analysis: Where We Win and Where We Don't}
% ============================================================

\subsection{Strengths}

\begin{itemize}
  \item \highlight{Full physics, zero approximation.} The production model
    runs the complete pvlib pipeline on AWS Lambda --- Perez transposition,
    SAPM thermal model, load-dependent inverter curve --- with no accuracy
    compromise from edge-side simplifications.
  \item \highlight{Serverless, scales to any fleet.} AWS Lambda charges
    per invocation only, with zero idle infrastructure and zero server
    provisioning or maintenance overhead (``NoOps''). The same physics
    engine serves one plant or one hundred --- monitoring cost scales
    linearly with fleet size, not with a per-kW subscription.
  \item \highlight{Multi-orientation aware.} Correctly handles 8 independent
    roof planes with different tilts and azimuths. Most commercial expected-gen
    tools use a single representative orientation.
  \item \highlight{Transparent and auditable.} Every parameter (losses, IAM
    table, thermal model, inverter curve) is stored in a human-readable JSON
    config and traceable to a PVsyst report. Nothing is a black box.
  \item \highlight{Cascade data resilience.} Four irradiance sources with
    per-field fallback means predictions keep running even when individual
    APIs fail.
  \item \highlight{Calibration included.} The calibration pipeline is
    fully automated. Adapting the engine to a new plant --- different
    module type, inverter, or geometry --- takes hours, not weeks.
  \item \highlight{Cost.} The entire stack uses free APIs (NASA POWER, ERA5,
    Open-Meteo) as primary sources. Total infrastructure cost can be zero.
\end{itemize}

\subsection{Weaknesses}

\begin{itemize}
  \item \highlight{No near-shading.} If a plant has significant row-to-row,
    tree, or structure shading that changes throughout the day, our model will
    over-predict. This is the single largest accuracy risk on constrained
    rooftops.
  \item \highlight{Irradiance-bounded accuracy.} On days with deep convective
    cloud cover or local monsoon, all reanalysis and satellite products
    diverge from ground truth simultaneously. No physics model can overcome
    incorrect inputs. This is the fundamental limit shared by every tool
    in the market on heavily overcast days. However, this boundary is actively
    shifting: AI-driven atmospheric foundation models (ECMWF AIFS, Google
    GraphCast) are demonstrably outperforming traditional numerical weather
    prediction --- and their open-weight versions are now publicly available
    for integration into irradiance cascades.
  \item \highlight{Static loss calibration.} Loss factors (soiling rate,
    mismatch) are set at commissioning and do not self-update. Over a
    multi-year horizon this introduces drift between the expected-generation
    baseline and actual plant behaviour.
\end{itemize}

% ============================================================
\section{Who This Is For}
% ============================================================

\subsection{Primary Audience}

\begin{itemize}
  \item \textbf{Solar O\&M engineers} managing commercial PV sites (50\,kW
    to 5\,MW) who need a physics-correct performance ratio baseline, not a
    nameplate-capacity average --- particularly those operating under
    \textbf{IEC\,61724-1:2021} performance assessment requirements that
    mandate auditable, physics-consistent expected yield.
  \item \textbf{IoT platform developers} building solar monitoring on AWS
    IoT Core, Azure IoT Hub, or similar --- who need a serverless expected-gen
    function that can be triggered per telemetry message.
  \item \textbf{Energy consultants and EPCs} who have PVsyst reports for a plant
    and want to translate those parameters into a live monitoring baseline without
    re-licensing PVsyst for every deployment.
  \item \textbf{DER aggregators and VPP operators} managing heterogeneous
    multi-vendor fleets for grid ancillary services or wholesale market bidding,
    where per-asset real-time expected generation must be computed at scale
    across diverse hardware without proprietary API dependency.
  \item \textbf{Research groups and universities} studying tropical solar
    performance who need a free, open, physics-correct baseline for comparison.
\end{itemize}

\subsection{Not the Right Tool For}

\begin{itemize}
  \item Bankable financial yield assessments (use PVsyst with certified
    irradiance data).
  \item String-level fault localisation (requires hardware-side diagnostics).
  \item Plants where near-shading \textgreater\ 5\% of annual yield (use a
    3-D shading engine first, then combine with our model).
  \item Real-time inverter MPPT optimisation (this is a monitoring / expected
    generation tool, not a control system).
\end{itemize}

\subsection{How to Adopt This Approach}

\begin{enumerate}
  \item \textbf{Extract plant parameters from your PVsyst commissioning report:}
    array orientations, module datasheet values, PVsyst loss factors, the IAM
    table, and the inverter efficiency curve.
  \item \textbf{Configure the physics engine} with those parameters in a
    structured plant configuration (location, geometry, losses, module and
    inverter specifications).
  \item \textbf{Connect to a free irradiance source} --- NASA POWER or
    Open-Meteo require no API key and cover any location globally.
  \item \textbf{Deploy the physics engine on AWS Lambda}, triggered by your
    IoT telemetry events. Each invocation returns expected AC power for the
    current irradiance and weather conditions.
  \item \textbf{Track performance ratio} as actual $\div$ expected generation
    on a 15-minute and daily basis. Persistent PR decline flags soiling or
    equipment degradation in near-real time.
\end{enumerate}

% ============================================================
\section{Why It Matters: The Value of Real-Time Accuracy}
% ============================================================

The 55\,kW plant in this study is a school building --- not a utility farm.
On a clear Sri Lankan day, it generates roughly 300--360\,kWh. A 10\%
soiling loss costs approximately 30--36\,kWh per day. That is a small number
in absolute terms, but it accumulates invisibly: without a physics-based
expected generation baseline, under-performance of this magnitude is
indistinguishable from a cloudy day.

\textbf{Four problems that a missing baseline makes structural:}

\begin{itemize}
  \item \textbf{Soiling goes undetected for weeks.} The only signal is a
    slow drift in daily output that looks like seasonal change. Globally,
    soiling reduces annual PV generation by 3--5\% on average and costs the
    solar industry an estimated \EUR{}4--7 billion per year \cite{skoplaki2009performance}.
    With a real-time performance ratio baseline, a 10\% soiling event appears
    within hours as a diverging PR, not at the next monthly report.
  \item \textbf{Performance ratio is meaningless without a physics model.}
    PR calculated as actual $\div$ nameplate rewards bad weather and penalises
    good weather. Only PR against \emph{physically modelled} expected
    generation --- driven by the same actual irradiance the plant saw ---
    separates equipment performance from meteorology.
  \item \textbf{Equipment faults are misread as cloud.} A degraded MPPT
    operating at 85\% efficiency produces an output signature identical to
    a partly cloudy afternoon, unless you have a model that accounts for
    both irradiance and expected electrical behaviour independently.
  \item \textbf{Operational decisions lack an evidence base.} Cleaning
    schedules, maintenance timing, and contractual yield reporting all require
    a transparent, auditable reference -- not a black-box dashboard number
    that cannot be queried or corrected.
\end{itemize}

\begin{insightbox}[The core principle]
The goal is not perfect forecasting. It is \emph{operational decision speed}:
detecting that something is wrong today, not next month, with a physics
explanation that an engineer can trace and verify. A continuously available
expected-generation baseline is the foundational layer that enables automated,
predictive O\&M --- shifting maintenance from reactive (dispatching a crew
after an inverter fully fails) to proactive (triggering a diagnostic the
moment the physics model diverges from actuals by a sustained threshold).
\end{insightbox}

\subsection{Data Ownership and Long-Term Value}

All predictions are stored locally in structured time-series files (hourly,
daily, monthly rolling). The entire system --- physics engine, configuration,
calibration pipeline --- is self-hosted. There is no vendor that can
revoke access, change API pricing, or sunset the product. Over years of
operation, accumulated expected-vs-actual records become an auditable asset
for O\&M reporting, degradation analysis, and project refinancing.

% ============================================================
\section{Conclusion \& Roadmap}
% ============================================================

What we built is a complete, physics-correct PV generation prediction
system: from a PVsyst commissioning report to a serverless AWS Lambda
function that fires at every telemetry tick, with no vendor lock-in,
no recurring API cost, and no accuracy compromise.

The validation on the 55\,kW school rooftop shows a system that correctly
separates weather-driven generation from equipment-driven generation ---
including the critical observation that during the December 2025 monsoon
period, model accuracy was limited not by the physics, but by the quality
of every available irradiance data source. That is a finding, not a failure.

What this system does not solve is the irradiance forecasting problem on
heavy cloud-cover days --- and that limitation is shared universally.

\textbf{Near-term roadmap:}
\begin{itemize}
  \item Formal annual energy parity validation against PVsyst simulation
  \item Automated soiling detection via rolling expected vs actual PR divergence
  \item Annual degradation rate tracking from multi-year PR trend analysis
  \item Physics-Informed Machine Learning (PIML) self-calibration layer:
    embedding physical solar geometry and thermal constraints directly into
    a self-updating parameter model, allowing loss factors to adapt to
    observed actuals without requiring full re-commissioning
  \item Multi-plant deployment with per-plant physics configuration
  \item Integration of next-generation AI weather models (e.g.\ ECMWF\,AIFS,
    Google\,GraphCast) into the irradiance cascade
  \item Open-source release with full documentation
\end{itemize}

\begin{insightbox}[The core thesis]
A physics engine embedded in your IoT pipeline --- not an annual report,
not a black-box dashboard, but a real-time serverless function that computes
\emph{what this plant should be generating right now} --- is the missing
layer in most solar SCADA deployments. The path from a PVsyst commissioning
report to a production AWS Lambda node can be measured in hours.
\end{insightbox}



\vspace{1em}

\noindent\rule{\linewidth}{0.4pt}

{\small\color{lightgrey}
\noindent\textit{For questions or collaboration on this work, connect with
S.\,M.\,Pasindu Upeksha Senevirathne via LinkedIn.}}

% ============================================================
%  BIBLIOGRAPHY
% ============================================================

\newpage
\printbibliography[title={References}]

\end{document}
